\documentclass[12pt]{article}
\usepackage[utf8]{inputenc}
\usepackage[T1]{fontenc}
\usepackage[margin=1in]{geometry}
\usepackage{setspace}
\usepackage{xcolor}
\usepackage[colorlinks=true, urlcolor=blue!60!black, linkcolor=blue!60!black]{hyperref}
\setlength{\parindent}{0em}
\setlength{\parskip}{0.6em}
\renewcommand{\baselinestretch}{1.0}
\begin{document}
\begin{center}
{\LARGE \textit{NFIB v. Sebelius}, 567 U.S. 519 (2012)\par}
\vspace{0.5em}
{\large Case Brief --- Two-Page Edition\par}
\end{center}

\section{Memory Jogger}
The ACA's individual mandate exceeds the commerce power but survives as a tax; threatening all Medicaid funds to force the expansion is unconstitutionally coercive (\href{https://supreme.justia.com/cases/federal/us/567/519/}{Justia}).

\section{Detailed Case Facts}
The Affordable Care Act required most Americans to carry insurance or pay a ``shared responsibility payment'' to the IRS, and required States to expand Medicaid to nearly all low-income adults or risk losing all federal Medicaid funds---over 10\% of most state budgets.

\section{Procedural History}
Twenty-six States, individuals, and NFIB sued; the Eleventh Circuit held the mandate beyond the commerce power but severable.

\section{Judicial Votes}
Roberts controlled every issue. Five Justices (Roberts and the Scalia--Kennedy--Thomas--Alito joint dissent) rejected the commerce theory; five (Roberts, Ginsburg, Breyer, Sotomayor, Kagan) upheld the mandate as a tax; seven found the Medicaid threat coercive; five limited the remedy.

\section{Holding}
An exaction that is modest, lacks scienter, is collected by the IRS, and leaves a lawful choice is a valid tax whatever its label. Threatening an existing grant that is a large share of state budgets to compel a new program is coercion.

\section{Analysis of Opinions}
\textbf{Roberts:} the power to regulate commerce presupposes existing activity, and a power to compel purchases would amount to a general police power. The tax reading was ``fairly possible,'' and the Medicaid threat was ``economic dragooning.'' \textbf{Ginsburg:} everyone enters the health-care market, so the mandate regulates economic activity and is necessary and proper. \textbf{Joint dissent:} calling a ``penalty'' a tax was ``verbal wizardry,'' and the whole Act should fall.

\section{Examples: Future Applications}
\textbf{Same-Side:} a modest IRS assessment on uninsured flood-zone homeowners is a valid tax. \textbf{Opposite-Side:} a criminally enforced mandate to buy electric vehicles is not. \textbf{Fence-Sitter:} a \$2,500 IRS assessment for skipping a \$50 vaccine may be prohibitory enough to become a penalty.

\section{Critique}
Barnett praised the commerce limits; Barrett said the tax reading strained the text. Progressive critics call the activity/inactivity line arbitrary and note the coverage gap the Medicaid holding created.

\section{Key Quotations}
\begin{quote}
``The Framers gave Congress the power to regulate commerce, not to compel it, and for over 200 years both our decisions and Congress's actions have reflected this understanding. There is no reason to depart from that understanding now.''
\end{quote}
\noindent --- Chief Justice Roberts, 567 U.S. at 555.

\end{document}
